\section{Plano de Teste de Penetração}

\subsection{Objetivos do Pentest}

O teste de penetração tem como principais objetivos:

\begin{itemize}
\item Identificar vulnerabilidades presentes nos sistemas e aplicações do
  Laboratório de Segurança;
\item Validar a exploração das falhas identificadas durante a varredura de
  vulnerabilidades;
\item Avaliar os impactos das vulnerabilidades sobre a confidencialidade,
  integridade e disponibilidade das informações;
\item Verificar a eficácia dos controles de segurança implementados;
\item Produzir evidências técnicas que auxiliem na correção das falhas
  identificadas;
\item Elaborar recomendações para mitigação dos riscos encontrados.
\end{itemize}

\subsection{Metodologia}

O Pentest será conduzido seguindo boas práticas reconhecidas pela indústria,
incluindo referências do OWASP Testing Guide, PTES (Penetration Testing
Execution Standard) e MITRE ATT&CK.

As etapas previstas são:

\begin{enumerate}
\item Planejamento e validação do escopo;
\item Reconhecimento e coleta de informações;
\item Mapeamento de ativos, páginas e serviços;
\item Identificação de vulnerabilidades;
\item Validação técnica dos achados;
\item Exploração controlada das vulnerabilidades;
\item Avaliação dos impactos e riscos;
\item Registro das evidências obtidas;
\item Elaboração do relatório técnico final;
\item Apresentação dos resultados e recomendações.
\end{enumerate}

\subsection{Ferramentas Utilizadas}

As principais ferramentas previstas para execução dos testes são:

\begin{itemize}
\item OWASP ZAP;
\item Navegadores web para validação manual dos achados;
\item Ferramentas de inspeção de tráfego HTTP;
\item Ferramentas auxiliares de análise e documentação.
\end{itemize}

\subsection{Vulnerabilidades Prioritárias}

Com base na varredura realizada anteriormente com o OWASP ZAP, foram
identificadas vulnerabilidades que serão priorizadas durante o Pentest:

\begin{itemize}
\item SQL Injection (CWE-89);
\item Cross Site Scripting Refletido -- XSS (CWE-79);
\item Path Traversal (CWE-22);
\item External Redirect (CWE-601);
\item Ausência de proteção Anti-CSRF (CWE-352);
\item Ausência de Content Security Policy (CWE-693);
\item Ausência de proteção contra Clickjacking (CWE-693).
\end{itemize}

Essas vulnerabilidades foram identificadas na aplicação
\texttt{http://testasp.vulnweb.com} e representam os principais pontos de
atenção durante a execução do teste.

\subsection{Registro das Atividades}

Todas as atividades realizadas deverão ser registradas em diário de execução,
contendo:

\begin{itemize}
\item Data e horário da atividade;
\item Ferramenta utilizada;
\item Ação executada;
\item Sistema ou funcionalidade analisada;
\item Resultado obtido;
\item Evidências coletadas;
\item Observações relevantes.
\end{itemize}

\subsection{Modelo de Registro de Atividades}

\begin{table}[H]
\centering
\small
\begin{tabularx}{\textwidth}{|p{2.2cm}|p{1.0cm}|p{2.0cm}|p{4.25cm}|X|}
\hline
\textbf{Data} &
\textbf{Hora} &
\textbf{Ferramenta} &
\textbf{Atividade} &
\textbf{Resultado} \\
\hline

DD/MM/AAAA &
00:00 &
OWASP ZAP &
Varredura de vulnerabilidades &
Vulnerabilidades identificadas e registradas. \\
\hline

\end{tabularx}
\caption{Modelo de registro das atividades executadas durante o Pentest.}
\label{tab:registro_atividades}
\end{table}

\subsection{Documentação das Vulnerabilidades}

Cada vulnerabilidade identificada deverá ser documentada contendo:

\begin{itemize}
\item Nome da vulnerabilidade;
\item Descrição técnica;
\item Sistema ou URL afetada;
\item Evidência da exploração;
\item Classificação de risco;
\item CWE e WASC associados;
\item Impacto potencial;
\item Recomendação de correção.
\end{itemize}

\subsection{Evidências}

As evidências deverão incluir:

\begin{itemize}
\item Capturas de tela;
\item Requisições e respostas HTTP;
\item Payloads utilizados;
\item Logs relevantes;
\item Resultados obtidos pelas ferramentas empregadas.
\end{itemize}

\subsection{Relatório Final}

Ao término do Pentest será elaborado um relatório consolidado contendo:

\begin{itemize}
\item Resumo executivo;
\item Escopo do teste;
\item Metodologia utilizada;
\item Vulnerabilidades encontradas;
\item Evidências técnicas;
\item Avaliação de riscos;
\item Recomendações de mitigação;
\item Conclusões do trabalho realizado.
\end{itemize}

\subsection{Aspectos Contratuais}

O contrato formal deverá estabelecer claramente:

\begin{itemize}
\item Escopo autorizado para os testes;
\item Responsabilidades das partes envolvidas;
\item Regras de comunicação;
\item Limitações de responsabilidade;
\item Requisitos de confidencialidade;
\item Direitos sobre a propriedade intelectual dos relatórios e
evidências produzidos.
\end{itemize}

Essas disposições deverão estar alinhadas ao Acordo de Não Divulgação (NDA) e
às regras definidas para a execução do Pentest no Laboratório de Segurança.
